\begin{qfdefinition}[Time value of money: future and present value of a single amount, simple and compound interest, compounding frequency, stated and effective rates, and conversion between rate quotations]\label{def:time-value-compounding}
\qffield{Statement} Money held today can be deposited or invested to earn interest, so an amount today is worth more than the same amount received later; this is the \emph{time value of money} \citep[sec.~7.1]{principles-of-finance-2026}. The simplest cash flow is a \emph{lump sum}, one amount at one date. Any stream of cash flows can be divided into lump sums, so the arithmetic of a single amount is the basis for valuing every other stream (sec.~7.1). A \emph{timeline} sets out the dates, with today, time 0, on the left, and is a good habit before any calculation (sec.~7.3).

\emph{Compounding} moves a value forward in time. With \emph{simple interest} only the original principal earns interest; with \emph{compound interest} the interest of earlier periods also earns interest, ``interest that is earned on interest'', so that $FV=PV(1+r)^n$ (sec.~7.2). \emph{Discounting} moves a value back: $PV=FV/(1+r)^n$, and the rate used is called the \emph{discount rate} (sec.~7.4).

A quoted annual rate means nothing until its compounding frequency is stated. At 10 percent a year, 100 grows to 110 with annual compounding, to $100\times1.05^2=110.25$ when 5 percent is credited every six months, and to 110.38 with quarterly compounding; ``the compounding frequency defines the units in which an interest rate is measured'', as kilometres and miles measure distance \citep[chap.~4, sec.~4.4, pp.~101--102]{s7-hulloptions}. A rate $R$ compounded $m$ times a year grows $A$ to $A(1+R/m)^{mn}$ in $n$ years, and to $Ae^{Rn}$ with continuous compounding; rates with different frequencies are converted by equating these growth factors (eqs.~(4.1)--(4.2), pp.~102--103). Hull measures rates with continuous compounding unless he states otherwise (p.~103).

OpenStax calls the quoted rate the \emph{stated} rate and the rate that includes compounding within the year the \emph{effective} rate: 1.5 percent a month is a stated 18 percent a year but an effective $1.015^{12}-1=19.562$ percent; with annual compounding the two coincide \citep[sec.~8.4]{principles-of-finance-2026}. Money-market instruments accrue simple interest in proportion to time, measured by a \emph{day count}: interest between two dates is the number of days between them divided by the days in the reference period, times the interest for that period; under actual/360, 90 days earn exactly a quarter of the quoted rate \citep[chap.~6, sec.~6.1, pp.~152--153]{s7-hulloptions}.

\qffield{Notation} $PV$ value at time 0, $FV$ value at time $n$, $r$ the rate per period, $n$ the number of periods; $(1+r)^n$ is the growth factor and $(1+r)^{-n}$ the \emph{discount factor}. Simple interest: $FV=PV(1+rn)$. A nominal (stated) annual rate $R_m$ compounded $m$ times a year has periodic rate $R_m/m$; its effective annual rate is $\mathrm{EAR}=(1+R_m/m)^m-1$ and its continuously compounded equivalent $R_c=m\ln(1+R_m/m)$. A simple rate $R_s$ for an accrual fraction $\tau$ (days over the day-count base) grows a deposit by $1+R_s\tau$.

Interpretation, derived here: every quotation states one growth factor over a stated period, so conversion means equating growth factors over the same calendar time $T$ in years: $(1+R_m/m)^{mT}=e^{R_cT}=(1+\mathrm{EAR})^T$, and for a single simple-interest period of $T$ years, $1+R_s\tau=(1+\mathrm{EAR})^T$. Solving the single-amount equation for the other unknowns: $r=(FV/PV)^{1/n}-1$ and $n=\ln(FV/PV)/\ln(1+r)$. The \emph{rule of 72} estimates the years to double as 72 divided by the rate in percent (OpenStax, sec.~7.4); exactly, $n=\ln2/\ln(1+r)$.

\qffield{Conditions} The formulas assume one known rate per period, unchanged over the horizon, and cash flows at the stated dates. For $r>0$ a future value rises and a present value falls as $n$ or $r$ rises; at $r=0$, $PV=FV$. The rate per period must match the period of compounding, and the number of periods must count those periods, not years. A simple rate describes growth over its own term only: a 90-day rate reinvested at the same rate for four successive periods compounds. Its accrual fraction depends on the day count (\hyperref[conv:us-rate-quotation]{U.S. day counts}), so the same quoted number under actual/360 and actual/365 is not the same rate. Inflation and risk change which rate is appropriate, not this arithmetic (\hyperref[def:required-return-components]{required-return components}).

The rule of 72 is an approximation, close for rates of a few percent and poorer as rates rise: at 6 percent it gives 12 years against an exact 11.8957, at 36 percent 2 years against 2.2542.

\qffield{Practice} Examples, original. \emph{Single amounts.} 1{,}000 at 5 percent a year for three years grows to $1{,}000\times1.05^3=1{,}157.625$ with compound interest and to 1{,}150 with simple interest; the difference, 7.625, is interest on interest. The present value of 1{,}000 due in three years at 5 percent is 863.8376. Growth from 1{,}000 to 1{,}500 in six years is $1.5^{1/6}-1=6.9913$ percent a year, and money doubles at 6 percent in $\ln2/\ln1.06=11.8957$ years.

\emph{Converting quotes.} A nominal 12 percent compounded monthly credits 1 percent a month and grows 100 to 112.6825 in a year. The same growth is an effective annual rate of 12.6825 percent, 12.3040 percent compounded semiannually ($2(1.01^6-1)$), 12.1204 percent compounded quarterly and 11.9404 percent continuously compounded ($12\ln1.01$). A 90-day deposit at 5 percent simple, actual/360, grows by $1+0.05\times90/360=1.0125$. If the same deposit could be renewed at that rate through a 365-day year, the effective annual rate would be $1.0125^{365/90}-1=5.1671$ percent, and the continuously compounded rate $\ln(1.0125)\times365/90=5.0380$ percent. Conversely, an effective 5 percent a year corresponds to a 90-day actual/360 simple rate of $(1.05^{90/365}-1)\times360/90=4.8412$ percent.

\qffield{Pitfalls} Multiplying a periodic rate by the number of periods in a year gives the stated rate, not the effective one (18 against 19.562 percent above). Rates quoted with different compounding frequencies or day counts must be converted before they are compared or combined. OpenStax describes the discount rate as ``the inverse of the growth rate'' (sec.~7.4); the rate is the same $r$ in both directions, and it is the discount factor $(1+r)^{-n}$ that is the reciprocal of the growth factor.

\qffield{See also} \hyperref[meth:cash-flow-streams]{valuing cash-flow streams}; \hyperref[def:interest-rate-measures]{interest-rate measures} (reference rates, zero and par yields); \hyperref[def:bond-yield-spread-measures]{bond yield measures} (periodicity conversions for bond yields); \hyperref[meth:floating-money-market-yields]{money-market yields} (discount-basis and money-market quotations).

\end{qfdefinition}

\begin{qfmethod}[Valuing dated cash-flow streams: discounting each amount, ordinary annuities and annuities due, perpetuities, growing perpetuities and growing annuities, and payments more frequent than annual]\label{meth:cash-flow-streams}
\qffield{Statement} A stream of unequal amounts, a \emph{mixed stream}, is valued by finding the present or future value of each amount and adding the results; the shortcuts for equal payments cannot be used \citep[sec.~9.1]{principles-of-finance-2026}. That values at a common date can be added is \hyperref[res:cash-flow-additivity]{cash-flow additivity}.

An \emph{annuity} is a stream of fixed periodic payments. In an \emph{ordinary annuity} the first payment is one period from now; in an \emph{annuity due} it is immediate, so its value is the ordinary annuity's multiplied by $(1+r)$ (sec.~8.2). OpenStax states that with other facts unchanged the annuity due is always larger (sec.~8.2). For positive payments that holds only when $r>0$; the two are equal when $r=0$, and the annuity due is smaller when $-1<r<0$. A \emph{perpetuity} is a series of payments that continues forever, such as a fixed preferred dividend or an endowment; when the first payment is one period away, its present value is the payment divided by the required rate of return, and a \emph{growing perpetuity} divides by the required rate minus the growth rate (sec.~8.1). For monthly or other non-annual payments, the annual rate is divided by the number of periods in a year and the number of years multiplied by it (sec.~8.3).

\qffield{Notation} Payment $C$, rate $r$ per payment period, $N$ payments, $v=1/(1+r)$. Interpretation, derived here by summing a geometric series:
\[ PV_{\mathrm{ord}}=C\sum_{t=1}^{N}v^t=C\,\frac{1-(1+r)^{-N}}{r},\qquad PV_{\mathrm{due}}=(1+r)\,PV_{\mathrm{ord}},\qquad FV_{\mathrm{ord}}=C\,\frac{(1+r)^N-1}{r}=(1+r)^N PV_{\mathrm{ord}}. \]
For $r>0$, $C/r$ is the time-0 value of a level perpetuity whose first payment is at $t=1$. The ordinary annuity is that perpetuity less a perpetuity whose first payment is at $N+1$, worth $(1+r)^{-N}C/r$ at time 0: $C/r-(1+r)^{-N}C/r$. With first payment $C_1$ at time 1 growing at $g$ per period, the \emph{growing annuity} and growing perpetuity are
\[ PV=\frac{C_1}{r-g}\Bigl[1-\Bigl(\frac{1+g}{1+r}\Bigr)^{N}\Bigr]\ (r\neq g),\qquad PV=\frac{NC_1}{1+r}\ (r=g),\qquad PV_{\infty}=\frac{C_1}{r-g}\ (r>g). \]

\qffield{Conditions} One rate for every period, payments at period ends unless the annuity is due, and payments that are certain or replaced by their expected values. In the growing formulas $C_1$ is the first payment, one period away; OpenStax applies $C/(R_s-G)$ to the current dividend, which treats that dividend as the next one (sec.~8.1; \hyperref[meth:instrument-present-value]{the $D_1$ pitfall}). \emph{Domains.} Every finite-stream formula needs $r>-1$, so that each discount factor $(1+r)^{-t}$ is defined and positive. The level-annuity closed forms also need $r\neq0$; at $r=0$ the sums give $PV_{\mathrm{ord}}=PV_{\mathrm{due}}=FV_{\mathrm{ord}}=NC$, the limit of the formulas as $r\to0$. The growing annuity needs $g>-1$ as well, so that payments keep their sign, and uses the $r=g$ form when the rates are equal. The level perpetuity $C/r$ needs $r>0$, and the growing perpetuity needs $-1<g<r$: only then is the ratio $(1+g)/(1+r)$ below 1. When $g\ge r$, or $r\le0$ for a level perpetuity, the terms do not shrink and the sum has no finite value.

Dividing an annual rate by $m$ is correct only for a nominal rate compounded $m$ times a year. An effective annual rate converts to the periodic rate $(1+\mathrm{EAR})^{1/m}-1$ (\hyperref[def:time-value-compounding]{rate conversion}).

\qffield{Practice} Examples, original. \emph{Mixed stream.} 300, 500 and 700 at years 1, 2 and 4, at 6 percent, are worth $300/1.06+500/1.06^2+700/1.06^4=1{,}282.4827$. \emph{Level annuity.} Five annual payments of 1{,}000 at 6 percent are worth 4{,}212.3638, by the formula and by discounting each payment; as an annuity due, 4{,}465.1056; their value at year 5 is 5{,}637.0930. At a zero rate both annuities are worth 5{,}000. At $-2$ percent the ordinary annuity is worth 5{,}314.5809 and the annuity due $0.98\times5{,}314.5809=5{,}208.2892$, which is less. \emph{Growth.} A perpetuity whose first payment of 100 arrives in a year and grows at 3 percent is worth $100/(0.08-0.03)=2{,}000$ at 8 percent; if 100 has just been paid, the next payment is 103 and the value 2{,}060. Ten such growing payments are worth 755.0134; with $g=r=8$ percent, ten payments starting at 100 are worth $10\times100/1.08=925.9259$, since each payment's growth cancels its discounting.

\emph{Monthly payments.} 200 a month for 36 months at a nominal 9 percent compounded monthly, 0.75 percent a month, is worth 6{,}289.3611. If instead 6 percent is an effective annual rate, the monthly rate is $1.06^{1/12}-1=0.48676$ percent and twelve monthly payments of 100 are worth 1{,}162.88; using 0.5 percent gives 1{,}161.89.

\qffield{Pitfalls} Applying an annuity formula to unequal payments. Treating an annuity due as ordinary, or the reverse, which misstates value by the factor $1+r$. Assuming that the annuity due is worth more at every rate, which fails when $r\le0$. Applying $C/r$ to a perpetuity whose first payment is immediate; its value is $C+C/r$. Using the payment just made in place of the next payment in a growing formula. Applying the growing-perpetuity formula when $g\ge r$. Adding amounts at different dates without first moving them to one date.

\qffield{See also} \hyperref[meth:loan-amortisation]{loan amortisation}; \hyperref[def:interest-rate-measures]{interest-rate measures} (annuity factors and their limits); \hyperref[meth:instrument-present-value]{present value of bond and equity cash flows}.

\end{qfmethod}

\begin{qfmethod}[Loan amortisation: the level payment from the annuity formula, the split of each payment into interest and principal, the schedule and its outstanding balance, and the total cost of borrowing]\label{meth:loan-amortisation}
\qffield{Statement} \emph{Loan amortisation} is a schedule of how and when a debt is repaid with interest. For a fixed-rate loan with level payments, the payment is found by solving the present value of an ordinary annuity for the payment; a monthly loan uses the annual rate divided by 12 and the number of years times 12 \citep[sec.~8.3]{principles-of-finance-2026}. OpenStax's three-year auto loan of 32{,}000 at 6 percent costs about 973.50 a month (0.5 percent a month over 36 months). Each payment contains interest and principal; as the debt falls, the interest part falls and the principal part rises, so the debt declines slowly at first and faster later. The interest paid in total is the lender's revenue, the cost of credit, 3{,}046.08 in that example (sec.~8.3, Table~8.4).

\qffield{Notation} Loan $L$, rate $r$ per payment period, $N$ payments, level payment $A=Lr/\bigl(1-(1+r)^{-N}\bigr)$. The schedule starts at $B_0=L$; in period $k$ the interest is $I_k=rB_{k-1}$, the principal repaid $P_k=A-I_k$, and the balance $B_k=B_{k-1}-P_k$, ending at $B_N=0$. Interpretation, derived here: the balance after $k$ payments is the present value of the payments still due, $B_k=A\bigl(1-(1+r)^{-(N-k)}\bigr)/r$, which equals the loan grown at $r$ less the payments grown at $r$, $L(1+r)^k-A\bigl((1+r)^k-1\bigr)/r$. Successive principal parts grow by the factor $1+r$, $P_{k+1}=(1+r)P_k$, and total interest is $NA-L$. These formulas need $r>-1$ and $r\neq0$. At a zero rate no interest accrues and each payment repays principal only. The limits as $r\to0$ are $A=L/N$ and $B_k=L(1-k/N)$.

\qffield{Conditions} A fixed rate, level payments at the ends of equal periods, and no fees, taxes or prepayment; lenders' charges and other extras are excluded, as in OpenStax's example. The loan's quoted 6 percent is a nominal rate compounded monthly, so its effective annual rate is 6.1678 percent. A payment rounded to the cent leaves a small residual balance that is settled in the final payment. Floating-rate loans, prepayment and the amortising and sinking-fund structures of bonds are treated elsewhere (\hyperref[def:bond-cash-flow-contingencies]{bond cash-flow structures}).

\qffield{Practice} Example, original, small enough to check by hand: 10{,}000 repaid in four annual payments at 8 percent. The payment is $10{,}000\times0.08/(1-1.08^{-4})=3{,}019.2080$.
\begin{itemize}
\item Year 1: interest 800.0000, principal 2{,}219.2080, balance 7{,}780.7920.
\item Year 2: interest 622.4634, principal 2{,}396.7447, balance 5{,}384.0473.
\item Year 3: interest 430.7238, principal 2{,}588.4843, balance 2{,}795.5630.
\item Year 4: interest 223.6450, principal 2{,}795.5630, balance 0.
\end{itemize}
After year 2, the two payments still due are worth $3{,}019.2080\times(1-1.08^{-2})/0.08=5{,}384.0473$, the scheduled balance. Each principal part is the previous one times 1.08, and total interest is $4\times3{,}019.2080-10{,}000=2{,}076.8322$. At a zero rate the same loan is repaid by four payments of 2{,}500. The schedule's balances are 7{,}500, 5{,}000, 2{,}500 and 0, which is $10{,}000(1-k/4)$.

\emph{OpenStax's auto loan, recomputed.} The 36-month annuity factor at 0.5 percent is 32.8710 and the exact payment 973.5020. In month 1, interest is $0.005\times32{,}000=160.00$ and principal 813.50. With the exact payment, total interest is 3{,}046.07; the table's 3{,}046.08 reflects rounding to the cent.

\qffield{Pitfalls} Dividing an effective annual rate by 12 to get a monthly loan rate. Assuming that a level payment repays equal amounts of principal: in the example the principal parts rise from 2{,}219 to 2{,}796. Estimating the balance part-way through as the loan less a pro rata share of principal; the balance is the present value of the remaining payments.

\qffield{See also} \hyperref[meth:cash-flow-streams]{annuities}; \hyperref[def:time-value-compounding]{stated and effective rates}; \hyperref[meth:internal-rate-of-return]{internal rate of return} (the loan rate as the IRR of the lender's cash flows).

\end{qfmethod}

\begin{qfmethod}[Internal rate of return: definition, solution by bracketing and Newton steps, when it exists and is unique, lending against borrowing cash flows, and what it does not measure]\label{meth:internal-rate-of-return}
\qffield{Statement} The \emph{internal rate of return} (IRR) is ``the discount rate that sets the present value of the cash inflows equal to the present value of the cash outflows'', the rate at which net present value is zero \citep[sec.~16.3]{principles-of-finance-2026}. OpenStax describes it as the rate at which the NPV profile crosses the horizontal axis (sec.~16.3). Crossing is not necessary: any rate with $\mathrm{NPV}=0$ is an IRR, including one at which the profile touches the axis and turns back. For a project whose IRR exceeds the cost of capital, OpenStax's rule is to accept it. The rule is reliable for a conventional investment, outlays followed only by inflows, which has a single IRR with NPV positive below it and negative above (Notation). Whether that IRR is positive is a separate matter: it is positive exactly when the undiscounted sum of the cash flows is positive. The IRR has no closed form in general; it is found by trial and error, which calculators automate (sec.~16.3).

OpenStax lists its limits (sec.~16.3). It ``may not be a single mathematical solution'' when negative cash flows occur in more than one period: a facility that costs 350, earns 950 a year later and then costs 620 to convert has two IRRs. It ignores differences in scale. And it is described as assuming that cash flows are reinvested at the IRR. Between mutually exclusive projects ``the rule is to take the project with the highest NPV'': a 25{,}000 machine with an IRR of 13.20 percent is preferred at a 9 percent cost of capital to a 16{,}000 machine with 14.10 percent, because its NPV is higher (sec.~16.5, Table~16.11).

\qffield{Notation} $\mathrm{NPV}(r)=\sum_{t=0}^{T}CF_t(1+r)^{-t}$ for $r>-1$; the IRR solves $\mathrm{NPV}(r)=0$. Interpretation, derived here.
\begin{itemize}
\item \emph{Bracketing (bisection).} Find $-1<a<b$ with $\mathrm{NPV}(a)$ and $\mathrm{NPV}(b)$ of opposite signs; evaluate the midpoint and keep the half in which the sign changes. Because NPV is continuous for $r>-1$, the bracket contains a root, and each step halves it. If $\mathrm{NPV}$ is zero at an end, that end is a root. Opposite signs certify a \emph{crossing} root, but bracketing applies only where such a bracket exists, and it cannot find every IRR. It misses a root at which NPV touches zero without changing sign, and two roots between ends of the same sign.
\item \emph{Newton steps.} $r_{k+1}=r_k-\mathrm{NPV}(r_k)/\mathrm{NPV}'(r_k)$ with $\mathrm{NPV}'(r)=-\sum_t tCF_t(1+r)^{-t-1}$. Convergence near a root is fast, but a poor start can leave the region $r>-1$ or reach a different root; where a sign-change bracket exists, a safe solver brackets first; where none exists, candidates from Newton steps or a scan of NPV are accepted only after the validation below. Near a multiple root convergence is only linear.
\item \emph{Validating a solution.} Accept a numerical IRR only if it lies in the permitted domain $r>-1$ and its NPV residual is small relative to the size of the cash flows, say $|\mathrm{NPV}(r)|\le10^{-9}\sum_t|CF_t|$. At a multiple root a small residual fixes the rate less precisely than at a simple root.
\item \emph{How many IRRs.} With $x=1/(1+r)$, NPV is a polynomial in $x>0$. By Descartes' rule of signs, the number of positive roots, counted with multiplicity, is at most the number of sign changes in $CF_0,\dots,CF_T$ (zeros skipped) and differs from it by an even number. A double root counts twice although it is one distinct IRR. The argument assumes equally spaced periods, so that $t$ is an integer and NPV a polynomial in $x$. A \emph{conventional} investment, outflows followed by inflows, has one change and so at most one IRR, a simple root. It has exactly one, because NPV has the sign of the first outflow as $r\to\infty$ and of the last inflow as $r\to-1$; NPV is positive at rates below it and negative above, which is what the decision rule needs. The IRR is positive exactly when the undiscounted sum, $\mathrm{NPV}(0)$, is positive, and negative when that sum is negative; the rule holds in either case.
\end{itemize}

\qffield{Conditions} The bound on the number of IRRs does not guarantee one. If NPV is positive at every rate there is no IRR, although the stream may still be worth having. The rule ``accept if the IRR exceeds the cost of capital'' applies only to a conventional investment, as described in Notation. For a borrowing-type stream, cash received first and paid later, NPV rises with the rate, so the acceptance rule reverses: borrowing is worthwhile when its IRR is \emph{below} the cost of capital. For a stream with several IRRs, or with a multiple root, neither rule applies, and the decision rests on the NPV at the cost of capital.

Interpretation: the IRR is computed from the cash flows alone. The reinvestment statement concerns realised wealth. Compounding the interim inflows to the horizon at the IRR reproduces growth of the outlay at the IRR; at any other reinvestment rate the realised return differs. The NPV at the cost of capital needs no such step. The same equation defines a bond's yield to maturity (\hyperref[meth:instrument-present-value]{present value of instrument cash flows}) and a portfolio's money-weighted return (\hyperref[def:return-measures]{return measures}). For cash flows on irregular dates, $t$ is measured in years, so that each flow is discounted for its own date. NPV is then not a polynomial in $x$, and the root-counting argument above does not apply as stated.

\qffield{Practice} Examples, original unless attributed. \emph{Solving.} The stream $-1{,}000$, 300, 400, 500 has an NPV of 58.8271 at 6 percent and an IRR of 8.8963 percent, found by bisection on $[0,1]$. From $r_0=0$, Newton steps give 7.6923, 8.8737, 8.8963 percent. \emph{Two IRRs} (OpenStax's facility, $-350$, $+950$, $-620$): NPV is $-20$ at 0 percent, $+13.9053$ at 30 percent and $-30$ at 100 percent. The IRRs, from the quadratic in $x$, are 9.1797 and 62.2488 percent; the project is worth having only at costs of capital between them. \emph{An IRR without a crossing.} The stream $-100$, $+220$, $-121$ has $\mathrm{NPV}(r)=-\bigl(11/(1+r)-10\bigr)^2$. NPV is never positive, and it is zero only at 10 percent, its sole distinct IRR. That IRR is a double root of $-100+220x-121x^2=-(11x-10)^2$, and the two sign changes count it twice. NPV is $-1$ at 0 percent and $-0.6944$ at 20 percent, so no bracket has ends of opposite signs, and bisection cannot start. The candidate 10 percent is confirmed by its domain and by a residual of zero. At an 8 percent cost of capital the IRR exceeds the hurdle, yet NPV is $-0.0343$: the stream is neither a conventional investment nor a conventional borrowing, and neither IRR rule applies.

\emph{No IRR.} $+100$, $-250$, $+200$: the quadratic $100-250x+200x^2$ has a negative discriminant, so NPV never reaches zero; its minimum, 21.875, is at 60 percent. \emph{Borrowing.} $+1{,}000$ now and $-1{,}100$ in a year has an IRR of 10 percent; with money costing 8 percent its NPV is $-18.52$. The rule ``accept if the IRR exceeds the hurdle'' would wrongly accept it. \emph{Scale} (OpenStax's machines): the IRRs are 14.0855 and 13.2040 percent, the NPVs at 9 percent 2{,}835.62 and 3{,}970.67, and the NPV profiles cross at 11.9265 percent, the IRR of the difference between the two streams. Below that rate the larger machine adds more value.

\qffield{Pitfalls} Reporting one of several IRRs as ``the'' return. Ranking mutually exclusive projects by IRR. Applying the investment rule to a borrowing-type stream, whose IRR must instead be below the cost of capital. Applying either IRR rule to a stream with several IRRs or a multiple root. Concluding that there is no IRR because no sign change was found. Trusting a numerical solution without checking that it lies in $r>-1$ and that its NPV residual is small.

\qffield{See also} \hyperref[meth:capital-project-evaluation]{capital project evaluation} (NPV against IRR in capital budgeting); \hyperref[meth:capital-allocation-principles]{capital allocation} (capital rationing and unequal lives); \hyperref[def:return-measures]{money-weighted return}; \hyperref[meth:loan-amortisation]{loan amortisation}.

\end{qfmethod}

\begin{qfdefinition}[Required-return components: real rate, expected inflation and risk premia; the Fisher relation, exact for realised inflation, approximate in its additive form]\label{def:required-return-components}
\qffield{Statement} A \emph{nominal} rate is quoted in money; a \emph{real} rate measures the growth of purchasing power. The Fisher relation, $1+i=(1+R)(1+h)$, links them, with $i$ the nominal rate, $R$ the real rate and $h$ expected inflation: a 6 percent real rate with 2 percent expected inflation requires a nominal rate of 8.12 percent \citep[sec.~7.4]{principles-of-finance-2026}. The same text also uses the subtraction ``real interest rate $=$ nominal interest rate $-$ inflation rate'' (sec.~3.4), which is the first-order approximation of the exact relation. The short-term government borrowing rate is ``often referred to as the risk-free rate of interest'', and lenders ``charge a premium to borrowers who are higher risk'' (sec.~3.4).

A nominal government yield is not just the real rate plus expected inflation. In the Federal Reserve staff's term-structure model, nominal yield $=$ real yield $+$ expected inflation $+$ inflation risk premium, and real yield $=$ expected average future real short rate $+$ real term premium; the two premia are ``extra compensations bond investors demand for bearing inflation risks and real interest rate risks, respectively'' \citep[eqs.~(1)--(2)]{inflation-compensation-decomposition-2019}. Inflation-indexed bonds trade at a yield above the true real yield by a liquidity premium, so their breakeven (inflation compensation) equals expected inflation plus the inflation risk premium minus that liquidity premium (eqs.~(3)--(4)). A borrower that may default pays a credit spread over the risk-free rate \citep[chap.~4, sec.~4.1, p.~98]{s7-hulloptions}.

\qffield{Notation} $R$ real rate, $h$ inflation, $i$ nominal rate, all for the same horizon and compounding. Interpretation: with \emph{realised} inflation $h$, or within one specified inflation scenario, $R=(1+i)/(1+h)-1$ is an identity, the realised (ex post) real return, and $i=R+h+Rh$ exactly. Approximation: $i\approx R+h$, with error $Rh$; with continuously compounded rates, $\ln(1+i)=\ln(1+R)+\ln(1+h)$ holds exactly. Interpretation, a decomposition by convention: the required yield on a bond or the required return on an equity claim $=$ nominal risk-free yield of matching maturity $+$ premia for the claim's risks (credit or default, liquidity, and for equity an equity risk premium). Premia stated as spreads add only as an approximation to products of growth factors.

\qffield{Conditions} When a nominal rate is agreed, future inflation is unknown, so there are two real rates: the \emph{ex ante} real rate expected when the loan is made, $i-E\pi$, and the \emph{ex post} real rate actually realised, $i-\pi$; the nominal rate ``can adjust only to expected inflation'', so the Fisher effect is written $i=r+E\pi$ \citep[chap.~5, sec.~5-3]{macroeconomics-2019}. These are the additive forms. \emph{Uncertain inflation} (Interpretation, derived here): deflating a nominal rate by expected inflation is not the expected real return, because the realised real return $(1+i)/(1+h)-1$ is a convex function of $h$, so its mean exceeds its value at the mean (Jensen's inequality). With a fixed nominal return of 10 percent and inflation of 0 or 20 percent with equal probability, deflating by mean inflation, 10 percent, gives 0 percent, while the mean realised real return is $\tfrac12\bigl(10+(1.10/1.20-1)\times100\bigr)=0.8333$ percent. A required nominal yield under uncertainty therefore also contains compensation for inflation risk, which is what the stochastic decomposition below separates out; the teaching calculation in Practice assumes a single known inflation rate. The premia in the Federal Reserve decomposition are model estimates from an affine no-arbitrage term-structure model of U.S.\ Treasury and TIPS yields, not observed prices; the 2019 note updates estimates published for 1990--2013 (Kim, Walsh and Wei, introduction and model sections). Interpretation: credit, liquidity and equity premia are compensation the market requires; a component cannot be read from one yield without a model or a matched comparison, such as a breakeven or a credit spread over the same-maturity government yield.

\qffield{Practice} Example, original, with inflation treated as known: a real rate of 1.5 percent and inflation of 2.5 percent give a nominal rate of $(1.015)(1.025)-1=4.0375$ percent; the additive approximation, 4.0 percent, misses $Rh=0.0375$ percentage points. Adding an inflation risk premium of 0.30 points gives a nominal risk-free yield of 4.3375 percent. A corporate bond with a credit spread of 1.20 and a liquidity premium of 0.30 points then needs 5.8375 percent, and an equity claim with an equity risk premium of 5.0 points needs 9.3375 percent. The premia here are illustrative, not estimates.

\qffield{Pitfalls} The breakeven inflation rate is not the expected inflation rate; it also contains the inflation risk premium and, with the opposite sign, the indexed bond's liquidity premium (eq.~(4)). Subtracting inflation from a nominal rate is an approximation; at high inflation the error is material: at 20 percent inflation and a 10 percent nominal rate, the exact real rate is $-8.33$ percent, the approximation $-10$ percent. Do not mix compounding conventions when adding components.

\qffield{See also} \hyperref[def:return-measures]{return measures} (real return); \hyperref[def:tips-breakevens]{inflation-linked bonds and breakevens}; \hyperref[mod:forward-rates-fra]{forward rates} (term-structure theories); \hyperref[meth:instrument-present-value]{present value of instrument cash flows}.

\end{qfdefinition}

\begin{qfmethod}[Present value of bond and equity cash flows, and the implied yield, required return and growth rate that equate price to present value]\label{meth:instrument-present-value}
\qffield{Statement} A bond promises fixed interest payments, \emph{coupons}, at regular dates and repays its \emph{face} (par) value at maturity; its \emph{yield to maturity} is the single discount rate that equates the promised payments to the price. A bond's price is the present value of its coupons plus the present value of its face value, discounted at the yield to maturity per period; for semiannual coupons the annual coupon and the yield, quoted as a nominal annual rate with semiannual compounding, are halved and the number of periods doubled. The yield cannot be isolated algebraically and is found with a financial calculator or spreadsheet; a bond priced above par has a coupon above its yield, one below par a coupon below its yield \citep[sec.~10.2]{principles-of-finance-2026}. A bond can equally be priced by discounting each cash flow at the zero (spot) rate for its own maturity; its yield is the single rate that reproduces that price \citep[chap.~4, secs.~4.5--4.6, pp.~104--105]{s7-hulloptions}.

A share is valued by the present value of its expected dividends. With a constant dividend $D$ forever and required return $r$, the value is the perpetuity $D/r$; with dividends growing at a constant rate $g$, it is $D_0(1+g)/(r-g)=D_1/(r-g)$, which ``requires the return rate $r$ to be greater than the growth rate $g$'' \citep[sec.~11.2]{principles-of-finance-2026}. The same equation, given a price, can be solved for the growth it implies (sec.~11.2, Example~2).

\qffield{Notation} $C$ coupon per period, $F$ face value, $y$ yield per period, $N$ periods, $z_t$ zero rate for maturity $t$, annual compounding and cash flows at period ends:
\[ P=\sum_{t=1}^{N}\frac{C}{(1+y)^t}+\frac{F}{(1+y)^N}=C\,\frac{1-(1+y)^{-N}}{y}+\frac{F}{(1+y)^N},\qquad P=\sum_{t=1}^{N}\frac{CF_t}{(1+z_t)^t}. \]
Quotation (Interpretation): halving applies to a yield quoted as a nominal annual rate compounded semiannually, $y_{\mathrm{nom}}$, whose half-year rate is $y_{\mathrm{nom}}/2$. A yield quoted as an \emph{effective} annual rate $y_{\mathrm{eff}}$ has half-year rate $(1+y_{\mathrm{eff}})^{1/2}-1$; the two describe the same bond when $(1+y_{\mathrm{nom}}/2)^2=1+y_{\mathrm{eff}}$ (\hyperref[def:interest-rate-measures]{compounding conversions}). Equity: $V_0=D/r$ (constant dividend, as for non-callable fixed-rate preferred stock); $V_0=D_1/(r-g)$, $r>g$; two-stage growth: $V_0=\sum_{t=1}^{n}D_t(1+r)^{-t}+(1+r)^{-n}D_{n+1}/(r-g_L)$, the last term the present value of the terminal value at $n$. Interpretation, derived by rearranging: implied required return $r=D_1/P_0+g$, dividend yield plus growth; implied growth $g=(rP_0-D_0)/(P_0+D_0)$.

\qffield{Conditions} The bond formula values the cash flows on a coupon date; pricing between coupon dates needs accrued interest and day counts (\hyperref[conv:us-rate-quotation]{quotation conventions}). A single yield is a summary of the zero curve, not a rate earned on each cash flow, and realising it requires reinvestment at that yield. The dividend models are valuation identities given their inputs: the expected dividends, the constant $g<r$ forever, and the required return, which itself contains the premia of \hyperref[def:required-return-components]{required-return components}. Value is very sensitive to $r-g$. Implied returns and growth are what the price requires under the model, not forecasts.

\qffield{Practice} Examples, original, annual compounding. A zero-coupon bond paying 1{,}000 in three years at 4 percent is worth 888.9964. A three-year 5 percent annual-coupon bond, face 100, is worth 102.7751 at a 4 percent yield (premium bond, coupon above yield); at zero rates of 3.0, 3.8 and 4.1 percent it is worth 102.5709, and its yield, 4.0723 percent, lies between the zero rates and closest to the three-year rate. Priced at 98.00 the same bond yields 5.7447 percent, found by bisection.

A preferred share paying a fixed 6 a year at a required 7.5 percent is worth 80.00. A share with $D_0=2.00$, $g=4$ percent and $r=9$ percent is worth $2.08/0.05=41.60$, and at that price $2.08/41.60+0.04=9$ percent. With 10 percent growth for two years and 4 percent after, the dividends are 2.20 and 2.42, the terminal value at year 2 is $2.42\times1.04/0.05=50.336$, and the value is 46.4220. At a price of 52 and $r=9$ percent, the implied constant growth is $(0.09\times52-2)/54=4.9630$ percent. Each figure was also obtained by summing the cash flows directly.

\qffield{Pitfalls} The constant-growth formula fails when $g\ge r$. Use the dividend expected next period, $D_1$, not the last one paid. Match the rate to the cash-flow period and compounding: cash flows of 5 in six months and 105 in a year are worth 100.2219 at a 10 percent \emph{effective} annual yield (half-year rate 4.8809 percent), the same as at a nominal 9.7618 percent compounded semiannually; halving the effective 10 percent instead gives 100.00, a mispricing. A higher yield means a lower price for given cash flows.

\qffield{See also} \hyperref[meth:cash-flow-streams]{annuities and perpetuities}; \hyperref[meth:internal-rate-of-return]{internal rate of return} (solving for the yield); \hyperref[def:interest-rate-measures]{interest-rate measures} (zero and par yields); \hyperref[res:cash-flow-additivity]{cash-flow additivity}; \hyperref[def:return-measures]{return measures} (the IRR).

\end{qfmethod}

\begin{qfresult}[Cash-flow additivity and no arbitrage: the value of a sum of cash flows is the sum of their values, applied to implied forward rates, a forward on a dividend-paying share, forward exchange rates and a replicated option]\label{res:cash-flow-additivity}
\qffield{Statement} Two portfolios with identical payoffs in every state must have identical values today; ``if this were not the case, an arbitrageur could buy the less expensive portfolio and sell the more expensive one'' \citep[chap.~11, sec.~11.4, p.~255]{s7-hulloptions}. Applied to cash flows dated in the future, value is additive: a stream is worth the sum of its dated cash flows, each discounted at the rate for its date, and a position assembled from simpler instruments is worth the sum of their values. Four consequences follow.

\begin{itemize}
\item \emph{Forward rates.} Investing for $T_2$ years, or for $T_1$ years and then at the forward rate from $T_1$ to $T_2$, must give the same amount; a borrower and lender at the zero rates can lock in the forward rate today: borrowing for one year and investing for two creates an outflow at year 1 and an inflow at year 2, a loan made between those dates at the forward rate, and borrowing for the longer and investing for the shorter term locks in forward borrowing \citep[chap.~4, sec.~4.8, p.~110]{s7-hulloptions}. Rolling short borrowings over at whatever rates prevail later is not a lock-in; it is a bet that rates will differ from the forwards (Business Snapshot~4.1, p.~111).
\item \emph{Forward on an asset with known income.} Buying the asset with borrowed money and selling it forward costs nothing today. If part of the loan is repaid by the known income, the amount owed at maturity is $(S_0-I)e^{rT}$, with $I$ the present value of the income, and this must be the forward price. If the forward is higher, the trade earns the difference; if lower, shorting the asset, lending the proceeds and buying forward does (chap.~5, sec.~5.5, pp.~130--131 and Table~5.3; \hyperref[res:cost-of-carry-forward-price]{cost of carry}).
\item \emph{Forward exchange rates.} Borrowing domestic currency, converting at spot, investing the foreign currency and selling the proceeds forward must cost nothing at maturity, which gives covered interest parity, $F_0=S_0e^{(r-r_f)T}$ with continuous compounding (chap.~5, sec.~5.10, pp.~137--140; \hyperref[res:cost-of-carry-forward-price]{cost of carry}).
\item \emph{Options.} When the underlying can move only to two prices, a position of $\Delta$ shares and a riskless bond can reproduce an option's payoff in both states, so the option is worth that position (chap.~13, sec.~13.1, pp.~288--290); and a call plus a zero-coupon bond paying $K$ has the same payoff as a put plus the share, so $c+Ke^{-rT}=p+S_0$ (sec.~11.4, eq.~(11.6), p.~256).
\end{itemize}

\qffield{Notation} Interpretation, the discrete-compounding forms used in the examples. Zero rates $z_1,z_2$ (annual): $(1+z_2)^2=(1+z_1)(1+f_{1,1})$, so $f_{1,1}=(1+z_2)^2/(1+z_1)-1$. A share at $S_0$ paying a known dividend $D$ at $t_D<T$, rate $r$ effective per year: $F_0=\bigl(S_0-D(1+r)^{-t_D}\bigr)(1+r)^T=S_0(1+r)^T-D(1+r)^{T-t_D}$, the dividend carried forward to $T$. Spot $S_0$ and forward $F_0$ in domestic currency per unit of foreign currency, domestic and foreign rates $i_d,i_f$ for the same period: $F_0=S_0(1+i_d)/(1+i_f)$. One period, share $S_0$ moving to $S_u$ or $S_d$, option payoffs $f_u,f_d$, rate $r$ for the period: $\Delta=(f_u-f_d)/(S_u-S_d)$, bond holding $B=(f_d-\Delta S_d)/(1+r)$ in present value, option value $\Delta S_0+B$; parity $c+K/(1+r)=p+S_0$.

\qffield{Conditions} No arbitrage among participants who can trade the component instruments without transaction costs, borrow and lend at the stated rates and, for parity, hold European options with the same strike and expiry on a share that pays no dividend before expiry (Hull, secs.~5.3 and~11.4). The forward argument needs an investment asset whose income is known in amount and date; for consumption assets only one bound holds (\hyperref[res:cost-of-carry-forward-price]{cost of carry}). Additivity values claims whose cash flows are fixed or replicable; it does not say what an unreplicable payoff is worth. The two-state option model is a teaching device: real prices take many values, and the replication must be rebalanced (Hull, chap.~13). Forward rates and forwards are prices that can be locked in, not forecasts.

\qffield{Practice} Examples, original. \emph{Forward rate.} With annual zero rates of 3.0 and 3.8 percent, the one-year rate one year forward is $1.038^2/1.03-1=4.6062$ percent. Buying a two-year zero that pays 1{,}000 for 928.1225 and borrowing that amount for one year at 3.0 percent gives dated cash flows of 0 now, $-955.9662$ at year 1 (repaying the loan) and $+1{,}000$ at year 2: a loan of 955.9662 from year 1 to year 2 at $1{,}000/955.9662-1=4.6062$ percent, fixed today. The reverse trades, selling the two-year zero short and lending the 928.1225 for one year, give $+955.9662$ at year 1 and $-1{,}000$ at year 2, borrowing locked at the same rate. Repaying the year-1 debt by a new one-year loan is a different position: its rate is unknown today, and at 8 percent it would cost 1{,}032.4434 at year 2.

\emph{Forward on a dividend-paying share.} A share at 50 pays a dividend of 1.00 in six months, and the rate is 4 percent a year, effective. The dividend is worth $1.00/1.04^{0.5}=0.9806$ today, so the one-year forward is $(50-0.9806)\times1.04=50.9802$. Replication: borrow 0.9806 for six months, repaid by the dividend, and 49.0194 for a year, buy the share and sell it forward; at the year end the share is delivered for 50.9802, exactly the debt. At a quoted forward of 51.50 this trade earns 0.5198. At 50.50, shorting the share, lending 49.0194 for a year and 0.9806 for six months to pay the dividend, and buying forward earns 0.4802. With the rate expressed as $\ln1.04$ continuously compounded, Hull's $(S_0-I)e^{rT}$ gives the same 50.9802.

\emph{Forward exchange rate.} At 1.1000 USD per EUR, a one-year USD rate of 4.0 percent and EUR rate of 2.5 percent give $1.1000\times1.04/1.025=1.116098$ USD per EUR, a forward premium of 160.98 pips. Borrowing USD 1{,}100, buying EUR 1{,}000, depositing to EUR 1{,}025 and selling that forward returns USD 1{,}144, the amount owed.

\emph{Option.} A share at 50 moves to 60 or 40, and the period rate is 5 percent. A call struck at 50 pays 10 or 0, so $\Delta=0.5$ and the bond holding is $-20/1.05=-19.0476$: buy half a share and borrow 19.0476, at a value of 5.9524. The position pays 10 or 0, as the call does. Parity gives the put struck at 50 as $5.9524+50/1.05-50=3.5714$, which a separate replication of the put confirms.

\qffield{Pitfalls} Additivity requires the same date and the same currency: amounts must be discounted to a common date before they are added. A forward exchange rate quoted the other way round (foreign per domestic) inverts the formula. The replicating $\Delta$ and bond change when the share price or time changes; a static position replicates only the one-period payoff.

\qffield{See also} \hyperref[mod:forward-rates-fra]{forward rates and FRAs}; \hyperref[res:cost-of-carry-forward-price]{cost-of-carry forward price}; \hyperref[res:synthetic-short-put-call]{synthetic short through options} (parity with dividends); \hyperref[meth:instrument-present-value]{present value of instrument cash flows}.

\end{qfresult}
